problem:  
Let \(M^n \subset S^{n+p}(1)\) be a compact minimal submanifold, and let \(\|B\|^2\) denote the squared norm of its second fundamental form.  
Simons’ classical gap theorem (1968) asserts that if \(\|B\|^2 \le n\), then \(M\) is either totally geodesic (\(\|B\|^2 = 0\)) or one of the *Clifford-type* minimal submanifolds (\(\|B\|^2 = n\)). Later refinements due to Chern–do Carmo–Kobayashi, Yau, and Peng–Terng revealed additional pinching gap phenomena for selected low codimensions.  

**Open Problem (Quantitative Universal Gap Conjecture):**  
For each \(n\), does there exist a constant \( \varepsilon_n > 0 \) such that any compact minimal submanifold  
\[
M^n \subset S^{n+p}(1)
\]
satisfying  
\[
n < \|B\|^2 < n + \varepsilon_n
\]
must (up to congruence) be one of the known *homogeneous minimal examples* in the unit sphere—namely, the Clifford minimal hypersurfaces, the Veronese surface (\(n=2, p=4\)), or their standard isometric generalizations determined by isotropy representations of compact symmetric spaces?  
In particular, is there a uniform gap
\[
\inf\{\|B\|^2(M) - n : M\text{ minimal, not congruent to a homogeneous example}\} > 0
\]
depending only on \(n\) (possibly also on \(p\))?  

Why is it a "good" problem:  
This formulation extends the Simons gap theorem into a conjectural *universal rigidity domain* for minimal submanifolds in spheres. It converts a known discrete rigidity threshold into a quantitative question about how “close” a minimal immersion can be, in curvature space, to the classical homogeneous models.  

It is “good” for several interlocking reasons:  

1. **Conceptual simplicity, profound depth:** The statement involves only the elementary invariant \(\|B\|^2\) and the ambient unit sphere, yet its resolution would classify an entire geometric genus of minimal submanifolds.  

2. **Connection to core phenomena:** It directly refines one of the most fundamental discoveries in global differential geometry—the Simons gap—and asks whether a new level of geometric rigidity exists universally beyond codimension one.  

3. **Bridging disciplines:** The problem necessarily engages tools from geometric analysis (Simons-type identities, stability operator spectra), representation theory of compact Lie groups (classification of isotropy irreducible homogeneous submanifolds), and potentially Ricci flow techniques or calibrated geometry.  

4. **Open and plausible:** The suggested universal bound is not currently known to exist; partial results depend on \(n,p\) and special symmetries. A positive resolution would constitute a cornerstone classification theorem; even disproving it would illuminate the landscape of minimal immersions near the known gaps.  

Thus this problem exemplifies a *“narrowly stated, deeply powerful”* question—an archetypal research-level problem connecting curvature pinching rigidity, minimal submanifold theory, and global geometric structure.